%% Chapter 2: Literature Survey

\chapter{LITERATURE SURVEY}

%% ─────────────────────────────────────────────────────────────────────────────
\section{OVERVIEW}

Multi-objective UAV path planning has attracted sustained research attention because no
single route can simultaneously minimise distance, energy, and flight risk---these
objectives are structurally in conflict. The dominant approach in the literature is to
generate a \textbf{Pareto-optimal front} of non-dominated routes, each representing a
different trade-off across the objective space, and then select one route from that front
as the mission route.

The critical question that determines the practical value of any such system is:
\textit{how is the final route selected from the Pareto front?} The three most directly
relevant papers---reviewed in detail below---all answer this question using
\textbf{front diversity metrics}: entropy, geometric spread, or crowding distance. These
metrics select routes that are well-distributed across the Pareto front, but they are
entirely blind to the delivery context. None of the three papers allow what the UAV is
carrying to influence which route is chosen.

This is the structural research gap this project addresses.

%% ─────────────────────────────────────────────────────────────────────────────
\section{REVIEW OF RELATED WORK}

\subsection{Chokkalingam and Ganesan (2026)}

Chokkalingam and Ganesan \cite{chokkalingam2026} present a multi-objective UAV path
planning framework that uses adaptive constraint handling combined with an
\textbf{entropy-based criterion} to select a final route from the generated Pareto front.
The entropy metric measures the information content of the objective distribution across
Pareto points: a point with high entropy is selected because it maximises the diversity
of objective trade-offs available to the operator.

The system is evaluated on custom simulation environments and demonstrates strong Pareto
front coverage compared to standard MOPSO baselines. The entropy criterion is shown to
select routes that span a wider range of objective combinations than random or
distance-ranked selection.

\textbf{Limitation.} The entropy criterion selects solutions based on front diversity
alone. Payload metadata---fragility class, delivery urgency, thermal sensitivity---is
never incorporated into the selection process. An operator carrying fragile biological
samples and an operator carrying emergency insulin would receive the same selected route
from the same Pareto front, because the entropy of the front is independent of what is
being transported.

\subsection{Datsko et al. (2024)}

Datsko et al. \cite{datsko2024} propose a diversity-driven Pareto selection mechanism for
UAV routing that optimises \textbf{geometric spread} of candidate solutions across the
multi-objective space. The key idea is to select Pareto points that maximise the physical
volume of the objective space that is ``covered'' by the candidate solution set, ensuring
the operator has access to a geometrically diverse set of route options.

The paper demonstrates improvements in Inverted Generational Distance (IGD) over
diversity-agnostic selection methods and shows that geometrically spread solutions
provide better coverage of the true Pareto front when the true front is not known in
advance.

\textbf{Limitation.} The geometric-spread selection mechanism is guided entirely by
solution diversity in objective space. The delivery context---payload type, time window,
priority class---has no influence on which route is chosen. The mechanism is designed to
serve human decision-making (giving the operator a diverse set of choices) rather than
autonomous selection aligned with mission requirements.

\subsection{Ntakolia et al. (2023)}

Ntakolia et al. \cite{ntakolia2023} apply a \textbf{crowding-distance-based} Pareto
point selection method within a multi-objective UAV path planner. Crowding distance,
originally introduced in NSGA-II, measures how isolated each Pareto point is from its
neighbours in objective space: points with high crowding distance are in less dense
regions of the front and are preferred to maintain spread.

The paper shows that crowding-distance selection produces well-spread Pareto fronts for
UAV routing problems and compares favourably against NSGA-II and $\epsilon$-constraint
methods on grid-world airspace environments.

\textbf{Limitation.} Crowding distance is a purely geometric measure of front density.
It cannot encode mission-specific requirements. Payload fragility, temperature
sensitivity, or urgency class cannot influence route selection at any stage. Like entropy
and geometric-spread methods, crowding distance maximises front diversity without any
awareness of what is actually being delivered.

%% ─────────────────────────────────────────────────────────────────────────────
\section{LITERATURE SURVEY SUMMARY TABLE}

Table~\ref{tab:litsurvey} summarises all three papers across key dimensions relevant
to the proposed work.

\begin{table}[ht!]
\centering
\caption{Summary of Reviewed Multi-Objective UAV Path Planning Literature}
\label{tab:litsurvey}
\renewcommand{\arraystretch}{1.3}
\resizebox{\textwidth}{!}{%
\begin{tabular}{p{3.5cm} p{1.8cm} p{3.2cm} p{3.2cm} p{3.5cm}}
\toprule
\textbf{Author(s) \& Year} &
\textbf{Journal / Venue} &
\textbf{Technique} &
\textbf{Pareto Selection Method} &
\textbf{Limitation / Research Gap} \\
\midrule
Chokkalingam \& Ganesan (2026) \cite{chokkalingam2026} &
RAIRO -- Oper.\ Res. &
Multi-objective UAV path planning with adaptive constraint handling &
Entropy-based criterion (front diversity) &
Payload metadata never incorporated; fragile samples and emergency medicine
receive identical route selection \\
\midrule
Datsko et al.\ (2024) \cite{datsko2024} &
IEEE RA-L &
Diversity-driven Pareto selection for UAV routing &
Geometric spread (maximise volume of objective space covered) &
Selection is guided entirely by solution diversity; delivery context has
no influence on route choice \\
\midrule
Ntakolia et al.\ (2023) \cite{ntakolia2023} &
J.\ Intell.\ Robot.\ Syst. &
Crowding-distance Pareto selection in multi-objective UAV planner &
Crowding distance (NSGA-II style) &
Crowding distance cannot encode urgency, fragility, or thermal sensitivity;
mission context is absent from selection \\
\bottomrule
\end{tabular}
}
\end{table}

%% ─────────────────────────────────────────────────────────────────────────────
\section{RESEARCH GAPS}

\subsection{Common Gap Across All Three Papers}

All three reviewed papers select Pareto-optimal routes based on \textbf{front diversity
metrics}---entropy, geometric spread, and crowding distance respectively. None of them
allow the delivery context (what payload is being carried) to influence which route is
selected. This is the structural gap the proposed system directly addresses.

The gap is not a minor implementation detail---it is a fundamental design assumption.
Entropy, geometric spread, and crowding distance are all defined over the Pareto front
itself. They have no interface through which external mission context can enter. Extending
these methods to incorporate payload information would require replacing the selection
criterion entirely, which is precisely what the proposed Tier~2 utility scoring mechanism
does.

\subsection{Specific Gaps Addressed by This Work}

\textbf{Gap 1 --- No payload-aware Pareto selection.} All existing UAV multi-objective
planners apply a fixed selection criterion independent of payload type. The proposed
Tier~2 introduces a formula-driven mapping from four payload tags directly to Pareto
point selection.

\textbf{Gap 2 --- No re-use across mission types.} Existing planners re-run the full
multi-objective search when the mission context changes. The proposed two-tier separation
pre-computes the Pareto front once (Tier~1) and allows instantaneous route selection for
any new payload tag (Tier~2), without re-running the planner.

\textbf{Gap 3 --- No algebraic bridge between delivery context and routing objectives.}
The mapping between a payload description and the routing trade-off it implies (e.g.,
fragility $\rightarrow$ minimise risk, urgency $\rightarrow$ minimise distance) has
never been formalised algebraically. The proposed weight-derivation formula provides this
bridge, making the mapping deterministic, transparent, and reproducible.
